\begin{longtable}{@{}l p{13cm}@{}}
  \caption{Incorrectly answered Praxis Reading questions: gold answer vs.\ the model's prediction and reasoning.}\label{tab:praxis_errors} \\
  \toprule
  \endfirsthead
  \toprule
  \endhead
  \multicolumn{2}{@{}l}{\textbf{q008}} \\
  \textbf{Passage}   & Use the following passage pair to answer questions 7 through 10. \newline \# Passage 1 \newline The variety of species among the jumping spiders appears almost infinite. The genus, or rather family of Epeira, is here characterized by many singular forms; some species have \newline pointed coriaceous shells, others enlarged and spiny tibiae. Every path in the forest is barricaded with the strong yellow web of a species belonging to the same division with the Epeira clavipes of Fabricius, which was formerly said \newline by Sloane to make, in the West Indies, webs so strong as to catch birds. A small and pretty kind of spider, with very long forelegs, and which appears to belong to an undescribed genus, lives as a parasite on almost every one of these \newline webs. I suppose it is too insignificant to be noticed by the great Epeira, and is therefore allowed to prey on the minute insects, which, adhering to the lines, would otherwise be wasted. When frightened, this little spider \newline either feigns death by extending its front legs, or suddenly drops from the web. \newline Source: Excerpt from Some Wonderful Spiders by Charles Darwin \newline \# Passage 2 \newline There are, among spiders, two forms of protective modification: the first, including all cases of protective resemblance to vegetable and inorganic things---that is, all modifications of color \newline or of color and form that tend to make them inconspicuous in their natural relations---I shall call direct protection. The second form, which I shall call indirect protection, includes two classes, the spiders which are specially pro- \newline tected themselves and those which mimic other creatures which are specially protected. \newline Spiders are specially protected when they become inedible through the acquisition of hard plates and sharp spines. The modification of form is frequently accompanied by conspicuous colors, which warn their enemies that they belong to an unpalatable class. \newline Source: Excerpt from Protective Resemblances in Spiders by Elizabeth G. Pechham    \\
  \textbf{Question}  & 8. With which of the following statements about the undescribed genus of spider that lives in the West Indies would the author of Passage 1 most likely agree? \newline a. It deserves to be named. \newline b. It is unimportant. \newline c. It is uniquely attractive. \newline d. It is an enemy of many animals. \newline e. It is cowardly.   \\
  \textbf{Gold}      & C       \\
  \textbf{Model}     & B  \\
  \textbf{Reasoning} & \textless{}think\textgreater{} \newline We need to answer question 8: "With which of the following statements about the undescribed genus of spider that lives in the West Indies would the author of Passage 1 most likely agree?" Options: a. It deserves to be named. b. It is unimportant. c. It is uniquely attractive. d. It is an enemy of many animals. e. It is cowardly. \newline We need to analyze Passage 1. The passage describes a small and pretty kind of spider, with very long forelegs, and which appears to belong to an undescribed genus, lives as a parasite on almost every one of these webs. I suppose it is too insignificant to be noticed by the great Epeira, and is therefore allowed to prey on the minute insects, which, adhering to the lines, would otherwise be wasted. When frightened, this little spider either feigns death by extending its front legs, or suddenly drops from the web. \newline The author says "I suppose it is too insignificant to be noticed by the great Epeira, and is therefore allowed to prey on the minute insects, which, adhering to the lines, would otherwise be wasted." So the author thinks it's too insignificant to be noticed. That suggests the author thinks it's unimportant (b). Also, the author says "I suppose it is too insignificant to be noticed by the great Epeira". So the author likely agrees that it is unimportant (b). Option a: "It deserves to be named." The author says it belongs to an undescribed genus, but doesn't explicitly say it deserves to be named. The author might think it's not important enough to be named. Option c: "It is uniquely attractive." The passage says "a small and pretty kind of spider", but not uniquely attractive. Option d: "It is an enemy of many animals." Not mentioned. Option e: "It is cowardly." The passage says when frightened, it feigns death or drops from the web. That could be seen as cowardly, but the author likely doesn't think it's cowardly; it's just a survival mechanism. The question asks which statement the author would most likely agree with. The phrase "too insignificant to be noticed" suggests the author thinks it's unimportant. So answer b. \newline But we need to check if any other option could be plausible. The author says "I suppose it is too insignificant to be noticed by the great Epeira". That directly supports "It is unimportant." So answer b. \newline Thus answer: B. \newline But the instruction says "Answer with the letter(s) only (e.g., A or AB). No explanation needed." So we output "B". \newline \textless{}/think\textgreater{} \newline B  \\
  \midrule
  \multicolumn{2}{@{}l}{\textbf{q013}} \\
  \textbf{Passage}   & Use the following passage to answer questions 13 through 15. \newline The account of this second journey of the Polos may be read in the wonderful book which Marco afterwards wrote to describe the won-ders of the world. They went from Lajazzo \newline through Turcomania, past Mount Ararat, where Marco heard tell that Noah's ark rested, and where he first heard also of the oil wells of Baku and the great inland sea of Caspian. Past Mosul and Bagdad they went, through Persia, where \newline brocades are woven and merchants bring caravan after caravan of treasures, to Hormuz, on the Persian Gulf, into which port put the ships from India, laden with spices, drugs, scented woods, and jewels, gold tissues and elephants' \newline teeth. Here they meant to take ship, but they desisted, perhaps because they feared to trust themselves to the flimsy nailless vessels in which the Arabs braved the dangers of the Indian Ocean. So they turned north again and \newline prepared to make the journey by land. They traversed the salt desert of Kerman, through Balk and Khorassan to Badakhshan, where there are horses bred from Alexander the Great's steed Bucephalus, and ruby mines and \newline lapis lazuli. It is a land of beautiful mountains and wide plains, of trout streams and good hunting, and here the brothers sojourned for nearly a year, for young Marco had fallen ill in the hot plains: a breath of mountain air blows \newline through the page in which he describes how amid the clean winds his health came back to him. \newline Source: Excerpt from A Short History of the World by H.G. Wells    \\
  \textbf{Question}  & 13. The passage suggests that Marco Polo was \newline a. generally easy to convince of fantastical things. \newline b. a man who possessed multiple extraordinary skills. \newline c. a coward compared to Arabian sailors. \newline d. homesick much of his life. \newline e. often in poor health because of his exhausting lifestyle.   \\
  \textbf{Gold}      & B       \\
  \textbf{Model}     & A  \\
  \textbf{Reasoning} & \textless{}think\textgreater{} \newline The user wants me to answer a multiple-choice question based on the provided passage. \newline **Question 13:** The passage suggests that Marco Polo was... \newline **Options:** \newline a. generally easy to convince of fantastical things. \newline b. a man who possessed multiple extraordinary skills. \newline c. a coward compared to Arabian sailors. \newline d. homesick much of his life. \newline e. often in poor health because of his exhausting lifestyle. \newline **Analysis of the passage:** \newline 1. **"They went from Lajazzo through Turcomania, past Mount Ararat... where Marco heard tell that Noah's ark rested, and where he first heard also of the oil wells of Baku and the great inland sea of Caspian."** -\textgreater{} Marco hears about Noah's Ark (fantastical/biblical) and oil wells (unusual for the time). This suggests he was open to hearing about wondrous things. \newline 2. **"Past Mosul and Bagdad they went, through Persia... to Hormuz... Here they meant to take ship, but they desisted, perhaps because they feared to trust themselves to the flimsy nailless vessels in which the Arabs braved the dangers of the Indian Ocean."** -\textgreater{} Marco feared the Arab ships. Does this make him a "coward"? The text says "perhaps because they feared". It doesn't explicitly label him a coward compared to them, just that he was cautious. \newline 3. **"So they turned north again and prepared to make the journey by land."** -\textgreater{} Decision to go by land instead of sea. \newline 4. **"They traversed the salt desert of Kerman... to Badakhshan... where there are horses bred from Alexander the Great's steed Bucephalus, and ruby mines and lapis lazuli."** -\textgreater{} Marco describes these things. He mentions Alexander's horse (fantastical/legendary). \newline 5. **"It is a land of beautiful mountains and wide plains, of trout streams and good hunting, and here the brothers sojourned for nearly a year, for young Marco had fallen ill in the hot plains: a breath of mountain air blows through the page in which he describes how amid the clean winds his health came back to him."** -\textgreater{} Marco fell ill in the plains and recovered in the mountains. This suggests he was in poor health at that specific time, but does it suggest he was *often* in poor health because of his exhausting lifestyle? The text mentions he fell ill once ("young Marco had fallen ill in the hot plains"). It doesn't provide evidence of a pattern of poor health due to lifestyle. \newline **Evaluating the options:** \newline * **a. generally easy to convince of fantastical things.** The passage mentions Marco hearing about Noah's Ark (biblical/fantastical) and Alexander's horse (legendary). He records these in his book ("The account of this second journey of the Polos may be read in the wonderful book which Marco afterwards wrote to describe the won-ders of the world"). The fact that he writes about Noah's Ark and Alexander's horse suggests he was open to or believed in these wondrous/fantastical accounts. This fits well. \newline * **b. a man who possessed multiple extraordinary skills.** The passage doesn't mention specific skills (writing, navigation, etc.) beyond traveling and writing a book. Not supported. \newline * **c. a coward compared to Arabian sailors.** The text says they "desisted, perhaps because they feared". It doesn't explicitly call him a coward or compare him negatively to the Arabs. It's speculative ("perhaps"). \newline * **d. homesick much of his life.** The passage doesn't mention homesickness. \newline * **e. often in poor health because of his exhausting lifestyle.** The text mentions he fell ill once in the plains and recovered in the mountains. "Often" is too strong. There's no mention of an exhausting lifestyle causing it generally. \newline **Conclusion:** Option **a** is the best fit. Marco records fantastical things like Noah's Ark and Alexander's horse, indicating he was open to or believed in such stories. \newline **Double check:** \newline - The passage starts by saying the account is in his book describing "wonders of the world". \newline - It mentions Noah's Ark (fantastical/biblical). \newline - It mentions Alexander's horse (legendary). \newline - He writes about these in his book. \newline - This implies he was open to or believed in these wondrous/fantastical accounts. \newline Therefore, **a** is the correct answer. \newline **Format:** The user wants the letter(s) only (e.g., A or AB). No explanation needed. \newline **Final Answer:** a \newline \textless{}/think\textgreater{} \newline a  \\
  \midrule
  \multicolumn{2}{@{}l}{\textbf{q042}} \\
  \textbf{Passage}   & Use the following passage to answer questions 41 and 42. \newline Science is knowledge; it is what we know. But mere knowledge is not science. For a bit of knowledge to become a part of science, its relation to other bits of knowledge must be found. \newline In botany, for example, bits of knowledge about plants do not make a science of botany. To have a science of botany, we must not only know about leaves, roots, flowers, seeds, etc., but we must know the relations of these parts and of all the parts of a plant to one another. In other words, in science, we must not only know, we must not only have knowledge, but we must \newline know the significance of the knowledge, must know its meaning. This is only another way of saying that we must have knowledge and know its relation to other knowledge. \newline Source: Excerpt from The Science of Human Nature by William Henry Pyle    \\
  \textbf{Question}  & 42. The passage is primarily concerned with \newline a. defining the term knowledge. \newline b. explaining how scientists must use knowledge. \newline c. emphasizing the importance of the whole over the part. \newline d. proving that most people who call themselves scientists are frauds. \newline e. showing how one becomes a scientist.   \\
  \textbf{Gold}      & B       \\
  \textbf{Model}     & C  \\
  \textbf{Reasoning} & \textless{}think\textgreater{} \newline We need to answer question 42. The passage is about science being knowledge but not just knowledge; it's about the relations between bits of knowledge. The passage says: "Science is knowledge; it is what we know. But mere knowledge is not science. For a bit of knowledge to become a part of science, its relation to other bits of knowledge must be found." Then it gives example in botany: bits of knowledge about plants do not make a science of botany; we must know the relations of these parts and of all the parts of a plant to one another. So the passage emphasizes that science is not just isolated bits of knowledge but the relations between them, i.e., the significance and meaning of knowledge in relation to other knowledge. \newline Thus the passage is primarily concerned with emphasizing the importance of the whole over the part? Or maybe "explaining how scientists must use knowledge"? Let's examine each answer: \newline a. defining the term knowledge. The passage does not define knowledge; it says science is knowledge but not just knowledge. It doesn't define knowledge. \newline b. explaining how scientists must use knowledge. The passage says we must know the significance of knowledge, its meaning, its relation to other knowledge. That could be interpreted as explaining how scientists must use knowledge (i.e., they must understand its relations). But the passage is more about what constitutes science, not specifically about how scientists use knowledge. \newline c. emphasizing the importance of the whole over the part. The passage says bits of knowledge about plants do not make a science; we must know the relations of these parts and of all the parts of a plant to one another. This suggests that the whole (the relations) is more important than the parts (bits of knowledge). So c seems plausible. \newline d. proving that most people who call themselves scientists are frauds. No, not mentioned. \newline e. showing how one becomes a scientist. Not really. \newline Thus the best answer is c. However, we need to check if the passage is primarily concerned with emphasizing the importance of the whole over the part. The passage says: "bits of knowledge about plants do not make a science of botany. To have a science of botany, we must not only know about leaves, roots, flowers, seeds, etc., but we must know the relations of these parts and of all the parts of a plant to one another." This is indeed about the whole (the relations) being more important than the parts (bits of knowledge). So answer c. \newline But let's double-check the wording: "The passage is primarily concerned with" - what is the main point? The passage argues that science is not just isolated facts but the relations between them. That is akin to emphasizing the whole over the part. So c is correct. \newline Thus answer: C. \newline But the instruction says "Answer with the letter(s) only (e.g., A or AB). No explanation needed." So we just output "C". \newline \textless{}/think\textgreater{} \newline C  \\
  \midrule
  \multicolumn{2}{@{}l}{\textbf{q008}} \\
  \textbf{Passage}   & Use the following passage pair to answer questions 7 through 10. \newline \# Passage 1 \newline The variety of species among the jumping spiders appears almost infinite. The genus, or rather family of Epeira, is here characterized by many singular forms; some species have \newline pointed coriaceous shells, others enlarged and spiny tibiae. Every path in the forest is barricaded with the strong yellow web of a species belonging to the same division with the Epeira clavipes of Fabricius, which was formerly said \newline by Sloane to make, in the West Indies, webs so strong as to catch birds. A small and pretty kind of spider, with very long forelegs, and which appears to belong to an undescribed genus, lives as a parasite on almost every one of these \newline webs. I suppose it is too insignificant to be noticed by the great Epeira, and is therefore allowed to prey on the minute insects, which, adhering to the lines, would otherwise be wasted. When frightened, this little spider \newline either feigns death by extending its front legs, or suddenly drops from the web. \newline Source: Excerpt from Some Wonderful Spiders by Charles Darwin \newline \# Passage 2 \newline There are, among spiders, two forms of protective modification: the first, including all cases of protective resemblance to vegetable and inorganic things---that is, all modifications of color \newline or of color and form that tend to make them inconspicuous in their natural relations---I shall call direct protection. The second form, which I shall call indirect protection, includes two classes, the spiders which are specially pro- \newline tected themselves and those which mimic other creatures which are specially protected. \newline Spiders are specially protected when they become inedible through the acquisition of hard plates and sharp spines. The modification of form is frequently accompanied by conspicuous colors, which warn their enemies that they belong to an unpalatable class. \newline Source: Excerpt from Protective Resemblances in Spiders by Elizabeth G. Pechham    \\
  \textbf{Question}  & 8. With which of the following statements about the undescribed genus of spider that lives in the West Indies would the author of Passage 1 most likely agree? \newline a. It deserves to be named. \newline b. It is unimportant. \newline c. It is uniquely attractive. \newline d. It is an enemy of many animals. \newline e. It is cowardly.   \\
  \textbf{Gold}      & C       \\
  \textbf{Model}     & B  \\
  \textbf{Reasoning} & \textless{}think\textgreater{} \newline Okay, let's see. The question is about what the author of Passage 1 would most likely agree with regarding the undescribed genus of spider in the West Indies. \newline First, I need to look at the relevant part of Passage 1. The passage mentions a "small and pretty kind of spider, with very long forelegs, and which appears to belong to an undescribed genus." The author says it's a parasite on the webs of other spiders, probably the Epeira. Then it says, "I suppose it is too insignificant to be noticed by the great Epeira, and is therefore allowed to prey on the minute insects..." \newline So the key points here are that the spider is small, pretty, has long forelegs, is a parasite on the webs, and is considered insignificant by the author. The author also mentions that it's allowed to stay because the Epeira doesn't notice it. \newline Looking at the options: \newline a. It deserves to be named. The passage says it "appears to belong to an undescribed genus," which means it hasn't been named yet. The author might think it's worth naming since it's a new genus, but the question is about what the author would agree with. The author's tone here is more about its role in the ecosystem rather than its taxonomic importance. \newline b. It is unimportant. The author uses "I suppose it is too insignificant," which directly relates to option b. The word "insignificant" here suggests the author thinks it's not important in the context of the Epeira's web. \newline c. It is uniquely attractive. The passage says "pretty," but "uniquely attractive" might be too strong. The author doesn't emphasize its beauty as a key point. \newline d. It is an enemy of many animals. The spider is a parasite on the webs, but the passage doesn't mention it being an enemy of many animals. It's more about its interaction with the Epeira and insects. \newline e. It is cowardly. The spider drops from the web when frightened, which could be seen as cowardly, but the author doesn't use that term. The behavior is described as a survival mechanism, not necessarily cowardice. \newline The best fit seems to be option b. The author's use of "insignificant" directly supports that the spider is considered unimportant in the context of the Epeira's web. Even though it's a new genus, the author's focus here is on its role and behavior, not its taxonomic significance. So the answer is B. \newline \textless{}/think\textgreater{} \newline B  \\
  \midrule
  \multicolumn{2}{@{}l}{\textbf{q011}} \\
  \textbf{Passage}   & Use the following passage to answer questions 11 and 12. \newline Frederick Douglass lived so long, and played so conspicuous a part on the world's stage, that it would be impossible, in a work of the size of this, to do more than touch upon the salient \newline features of his career, to suggest the respects in which he influenced the course of events in his lifetime, and to epitomize for the readers of another generation the judgment of his contemporaries as to his genius and his character. \newline Douglass's fame as an orator has long been secure. His position as the champion of an oppressed race, and at the same time an example of its possibilities, was, in his own generation, as picturesque as it was unique; and his \newline life may serve for all time as an incentive to aspiring souls who would fight the battles and win the love of mankind. \newline Source: Frederick Douglass by Charles Chestnutt    \\
  \textbf{Question}  & 11. The passage suggests that Frederick Douglass \newline a. was most concerned with being loved by other people. \newline b. mostly influenced matters in his hometown. \newline c. had a uniquely engaging personality. \newline d. has been the topic of many long history books. \newline e. was mostly famous for being a great speaker.   \\
  \textbf{Gold}      & C       \\
  \textbf{Model}     & E  \\
  \textbf{Reasoning} & \textless{}think\textgreater{} \newline Okay, let's see. The question is asking what the passage suggests about Frederick Douglass. The options are a through e. \newline First, I need to look at the passage again. The first part says he lived so long and had such a prominent role that it's hard to cover all his career features in a work of this size. So the author is saying there's a lot to his career, but they can only touch on the main points. That might relate to option d, which says he's been the topic of many long history books. But the passage is talking about the current work being limited in size, not necessarily that there are many books about him. Hmm. \newline Next, the passage mentions his fame as an orator is secure. So that's a direct statement about his speaking skills. Option e says he was mostly famous for being a great speaker. The word "mostly" here is important. The passage says his fame as an orator is secure, but also talks about his role as a champion of an oppressed race and an example of its possibilities. So maybe he's known for more than just speaking, but the passage does highlight his oratory skills. However, the question is about what the passage suggests. The first sentence says it's impossible to cover all his career features, but the second sentence starts with his fame as an orator. So maybe e is correct, but I need to check other options. \newline Option c says he had a uniquely engaging personality. The passage says his position was "as picturesque as it was unique" and that his life serves as an incentive. That might imply his personality was engaging. But the passage doesn't use the word "personality" directly. It's more about his role and influence. \newline Option a says he was most concerned with being loved by others. The passage mentions his life as an incentive to "win the love of mankind," but that's in the context of fighting for his race, not necessarily his personal concern. So a might not be right. \newline Option b says he influenced matters in his hometown. The passage says he influenced the course of events in his lifetime, but not specifically his hometown. So b is probably incorrect. \newline So between c and e. The passage says his fame as an orator is secure, which supports e. But the "as picturesque as it was unique" part might support c. However, the question is about what the passage suggests. The first part says he played a conspicuous part on the world's stage, which implies a broader influence than just his hometown. The mention of his oratory fame is a specific point, but the passage also talks about his role as a champion. So maybe e is correct. Let me check again. The first sentence says it's impossible to cover all his career features, so the author is acknowledging his long and significant career. The second sentence starts with his fame as an orator, which is a key point. So the passage suggests that he was famous for being a great speaker. Therefore, the answer is e. \newline \textless{}/think\textgreater{} \newline E  \\
  \midrule
  \multicolumn{2}{@{}l}{\textbf{q013}} \\
  \textbf{Passage}   & Use the following passage to answer questions 13 through 15. \newline The account of this second journey of the Polos may be read in the wonderful book which Marco afterwards wrote to describe the won-ders of the world. They went from Lajazzo \newline through Turcomania, past Mount Ararat, where Marco heard tell that Noah's ark rested, and where he first heard also of the oil wells of Baku and the great inland sea of Caspian. Past Mosul and Bagdad they went, through Persia, where \newline brocades are woven and merchants bring caravan after caravan of treasures, to Hormuz, on the Persian Gulf, into which port put the ships from India, laden with spices, drugs, scented woods, and jewels, gold tissues and elephants' \newline teeth. Here they meant to take ship, but they desisted, perhaps because they feared to trust themselves to the flimsy nailless vessels in which the Arabs braved the dangers of the Indian Ocean. So they turned north again and \newline prepared to make the journey by land. They traversed the salt desert of Kerman, through Balk and Khorassan to Badakhshan, where there are horses bred from Alexander the Great's steed Bucephalus, and ruby mines and \newline lapis lazuli. It is a land of beautiful mountains and wide plains, of trout streams and good hunting, and here the brothers sojourned for nearly a year, for young Marco had fallen ill in the hot plains: a breath of mountain air blows \newline through the page in which he describes how amid the clean winds his health came back to him. \newline Source: Excerpt from A Short History of the World by H.G. Wells    \\
  \textbf{Question}  & 13. The passage suggests that Marco Polo was \newline a. generally easy to convince of fantastical things. \newline b. a man who possessed multiple extraordinary skills. \newline c. a coward compared to Arabian sailors. \newline d. homesick much of his life. \newline e. often in poor health because of his exhausting lifestyle.   \\
  \textbf{Gold}      & B       \\
  \textbf{Model}     & E  \\
  \textbf{Reasoning} & \textless{}think\textgreater{} \newline Okay, let's see. The question is asking what the passage suggests about Marco Polo. The options are a through e. \newline Looking at the passage, it mentions Marco hearing about Noah's ark resting on Mount Ararat and the oil wells of Baku and the Caspian Sea. Then later, when they get to Hormuz, they were going to take a ship but decided not to because they were afraid of the flimsy Arab ships. Then they went north again, and Marco fell ill in the hot plains, so they stayed in Badakhshan for a year where the mountain air helped him recover. \newline So, the key points here are that Marco was told about these amazing places and things, and he was concerned about the safety of the ships, which might indicate he was cautious. Also, he got sick in the hot plains and needed time to recover in the mountains. The question is about what the passage suggests about him. \newline Option e says he was often in poor health because of his exhausting lifestyle. The passage mentions he fell ill once, but it doesn't say he was often in poor health. It just says he got sick once and then recovered. So maybe e isn't the right answer. \newline Option a says he was generally easy to convince of fantastical things. The passage says he heard tell of Noah's ark and other things, but it doesn't say he believed them easily. It's more about what he heard, not his belief in them. \newline Option b says he had multiple extraordinary skills. The passage doesn't mention any skills, just his journey and illness. \newline Option c says he was a coward compared to Arabian sailors. The passage says they desisted from taking the Arab ships because they were flimsy, but it's not clear if Marco was the one who was afraid. The brothers might have made the decision together. So maybe c is not correct. \newline Option d says he was homesick much of his life. The passage doesn't mention homesickness at all. They stayed in Badakhshan for a year because of his health, not because he was homesick. \newline So the only one that seems to fit is e, but the passage only mentions one instance of illness. However, the question is about what the passage suggests. The fact that he fell ill in the hot plains and needed to stay in a cooler place might imply that his health was affected by the journey's conditions. But the option says "often in poor health because of his exhausting lifestyle." The passage doesn't state that he was often ill, just once. So maybe e is not correct. Wait, but the question is about what the passage suggests. The passage does mention that he fell ill in the hot plains, which could be part of an exhausting lifestyle. However, the answer might be e if the implication is that his health was often poor. But the passage only mentions one illness. Hmm. Maybe the answer is e, but I'm not sure. Let me check again. \newline The passage says "young Marco had fallen ill in the hot plains" and they stayed in Badakhshan for a year. So it's a specific instance. The option e says "often in poor health because of his exhausting lifestyle." The passage doesn't say "often," just once. So maybe e is incorrect. Then, perhaps the answer is a? Because he heard about these things. But the question is about what the passage suggests about him. The passage doesn't say he was easy to convince, just that he heard about them. So maybe none of the options are correct. Wait, but the options are given. Let me think again. \newline The passage mentions that Marco heard of oil wells and the Caspian Sea, and later when he was ill, the mountain air helped. The fact that he was ill in the hot plains might suggest that his health was affected by the journey's conditions. So maybe e is the answer. The passage does say he fell ill, which could be due to the exhausting journey. But the option says "often in poor health," which is not stated. The passage only mentions one time. So maybe e is not correct. Then, perhaps the answer is a? Because he heard of these things, but the question is about what the passage suggests about him. The passage doesn't say he was convinced of them, just that he heard about them. So maybe the answer is e, but I'm not sure. Wait, the question is about the passage suggesting something. The passage says he fell ill in the hot plains, which might indicate that his health was affected by the journey. So maybe e is the answer. But the passage doesn't say he was often in poor health. Hmm. Maybe the answer is e, but I'm not certain. Alternatively, maybe the answer is a. The passage mentions he heard of these things, but it doesn't say he was convinced. So maybe a is not correct. \newline Wait, the question is about what the passage suggests. The fact that he wrote about the wonders of the world in his book might imply he was open to hearing about fantastical things. So maybe a is correct. But the passage doesn't say he was easy to convince, just that he heard about them. So I'm a bit confused. Let me check the options again. The answer is probably e, but I'm not sure. Alternatively, maybe the answer is a. The passage says he heard of these things, but the question is about what the passage suggests about him. If he wrote a book about the wonders, maybe he was easy to convince. But the passage doesn't explicitly say that. Hmm. I think the answer is e, but I'm not 100\% sure. Wait, the passage says "young Marco had fallen ill in the hot plains," which might indicate that his health was affected by the journey. So maybe e is the answer. But the passage doesn't say he was often in poor health. So perhaps the answer is e, but the wording is tricky. The question is about what the passage suggests. The passage does mention one instance of illness, which could be due to the exhausting journey. So maybe e is the answer. I'll go with e. \newline \textless{}/think\textgreater{} \newline e  \\
  \midrule
  \multicolumn{2}{@{}l}{\textbf{q017}} \\
  \textbf{Passage}   & Use the following passage pair to answer questions 16 through 18. \newline \# Passage 1 \newline American independence, the beginnings of which we have just been considering, was accomplished after a long struggle. Many brave men fought on the battle-field, and many who never shouldered a musket or drew a sword exerted a powerful influence for the good of the patriot cause. One of these men was Benjamin Franklin. \newline He was born in Boston in 1706, the fifteenth child in a family of seventeen children. His father was a candle-maker and soap-boiler. Intending to make a clergyman of Benjamin, he sent him, at eight years of age, to a grammar-school, with the purpose of fitting him for college. The boy made rapid progress, but before the end of his first school-year his father took him out on account of the expense, and put him into a school where he would learn more practical subjects, such as writing and arithmetic. The last study proved very difficult for him. \newline Source: Excerpt from American Leaders and Heroes: A Preliminary Text-Book in United States History by Wilbur F. Gordy \newline \# Passage 2 \newline At ten years old I was taken home to assist my father in his business, which was that of a tallow-chandler and sope-boiler; a business he was not bred to, but had assumed on his arrival in New England, and on finding his dying trade would not maintain his family, being in little request. Accordingly, I was employed in cutting wick for the candles, filling the dipping mold and the molds for cast candles, attending the shop, going of errands, etc. \newline I disliked the trade, and had a strong inclination for the sea, but my father declared \newline against it; however, living near the water, I was much in and about it, learnt early to swim well, \newline and to manage boats; and when in a boat or canoe with other boys, I was commonly allowed to govern, especially in any case of difficulty; and upon other occasions I was generally a leader among the boys, and sometimes led \newline them into scrapes, of which I will mention one instance, as it shows an early projecting public spirit, tho' not then justly conducted. \newline Source: Excerpt from Benjamin Franklin: His Autobiography by Benjamin Franklin    \\
  \textbf{Question}  & 17. Which of the following is a difference between Passage 1 and Passage 2? \newline a. Passage 1 is more humorous than Passage 2. \newline b. Passage 1 tells the story of one person's life. \newline c. Passage 1 is about the life of an important American. \newline d. Passage 1 is more specific than Passage 2. \newline e. Passage 1 focuses more on a particular period of time.   \\
  \textbf{Gold}      & D       \\
  \textbf{Model}     & E  \\
  \textbf{Reasoning} & \textless{}think\textgreater{} \newline Okay, let's see. The question is asking for the difference between Passage 1 and Passage 2. The options are a through e. \newline First, I need to recall what each passage is about. Passage 1 is from a textbook by Wilbur F. Gordy, talking about Benjamin Franklin's role in American independence. It mentions his early education and how his father tried to make him a clergyman but switched to practical subjects. \newline Passage 2 is from Franklin's own autobiography. It describes his childhood work in his father's tallow-chandler and soap-boiler business, his dislike for the trade, and his early interest in the sea and leadership among boys. \newline Looking at the options: \newline a. Passage 1 is more humorous than Passage 2. Hmm, I don't see any humor in Passage 1. The autobiography might have more personal anecdotes, but the question is about the difference. Not sure yet. \newline b. Passage 1 tells the story of one person's life. Wait, Passage 1 is about Franklin's early life and his role in the patriot cause. Passage 2 is also about Franklin's early life. But the textbook passage is part of a larger work about American leaders and history, so maybe it's not the whole story of his life. The autobiography would be more about his entire life. So maybe this isn't the difference. \newline c. Passage 1 is about the life of an important American. Both passages are about Franklin, who is an important American. So this might not be a difference. \newline d. Passage 1 is more specific than Passage 2. Let me check. Passage 1 says he was born in Boston in 1706, 15th of 17 children. His father was a candle-maker and soap-boiler. Passage 2 gives more details about his work in the business, like cutting wick, filling molds, etc. So Passage 2 is more specific. Therefore, Passage 1 is less specific. So option d is incorrect. \newline e. Passage 1 focuses more on a particular period of time. Passage 1 talks about his early education and the context of American independence. Passage 2 is about his childhood work and early interests. Both are about different periods? Wait, Passage 1 mentions his father sending him to grammar school at 8, then taking him out. Passage 2 starts when he's 10, working in the business. So Passage 1 covers up to when he was 8, and Passage 2 from 10 onwards. But the question is about the difference in focus. Passage 1 is part of a history textbook, so it's discussing his role in the patriot cause and his early life in that context. Passage 2 is from his autobiography, giving personal details about his childhood. So maybe Passage 1 is more about the broader historical context and his role, while Passage 2 is about his personal experiences. But the options don't mention that. Let me check the options again. \newline Option e says Passage 1 focuses more on a particular period. But Passage 1 is about his early education and the start of American independence. Passage 2 is about his work in the family business and early interests. So Passage 1 is about his early education (age 8), and Passage 2 is about age 10 and his work. So Passage 1 is more about a specific time (early education), while Passage 2 is about another specific time (childhood work). But the question is about the difference between the two passages. \newline Looking back at the options, the correct answer is probably e. Because Passage 1 is part of a history textbook focusing on his role in the patriot cause and his early education, while Passage 2 is from his autobiography, focusing on his childhood work and personal inclinations. So Passage 1 is more about the broader historical context and his early life in relation to that, whereas Passage 2 is about a specific period of his childhood. Wait, but the question is which is a difference between the two. Let me check the options again. \newline Option e says Passage 1 focuses more on a particular period. But Passage 1 is part of a textbook that's about American independence and includes his early life as part of that. Passage 2 is a personal account of his childhood. So Passage 1 is part of a larger narrative about American history, while Passage 2 is a personal narrative. The difference might be that Passage 1 is part of a story about one person's life (Franklin's) in the context of American history, while Passage 2 is also about his life but more personal. \newline Wait, option b says Passage 1 tells the story of one person's life. But Passage 1 is part of a textbook about American leaders and heroes, so it's not the whole story of his life, just a part. Passage 2 is from his autobiography, which would be more about his entire life. So maybe option b is incorrect. \newline Option e: Passage 1 focuses more on a particular period. The first passage is about his early education and how his father's decisions affected his path, which is a specific time in his life. The second passage is about his work in the family business and his early interests, another specific period. So the difference is that Passage 1 is about a particular period (early education), while Passage 2 is about another period (childhood work). Therefore, option e is correct. \newline But wait, the question is which is a difference between Passage 1 and Passage 2. Let me check the options again. \newline Option e: Passage 1 focuses more on a particular period of time. The first passage is about his early education and the start of American independence. The second is about his childhood work. So Passage 1 is focusing on a specific time (early education) in the context of the broader history of American independence. Passage 2 is focusing on his work in the family business. So the difference is that Passage 1 is part of a history textbook discussing his role in the patriot cause and his early life in that context, while Passage 2 is a personal account of his childhood work. Therefore, the correct answer is e. \newline \textless{}/think\textgreater{} \newline e  \\
  \midrule
  \multicolumn{2}{@{}l}{\textbf{q024}} \\
  \textbf{Passage}   & Use the following passage to answer questions 24 and 25. \newline A history of architecture is a record of man's efforts to build beautifully. The erection of structures devoid of beauty is mere building, a trade and not an art. Edifices in which \newline strength and stability alone are sought, and in designing which only utilitarian considerations have been followed, are properly works of engineering. Only when the idea of beauty is added to that of use does a structure take its place \newline among works of architecture. We may, then, define architecture as the art which seeks to harmonize in a building the requirements of utility and of beauty. It is the most useful of the fine arts and the noblest of the useful arts. \newline Source: Excerpt from A Text-Book of the History of Architecture by A.D.F. Hamlin    \\
  \textbf{Question}  & 24. The final sentence of the passage makes which of the following errors in logic? \newline a. incomplete comparison \newline b. faulty generalization \newline c. circular reasoning \newline d. fallacy of composition \newline e. moral equivalence   \\
  \textbf{Gold}      & B       \\
  \textbf{Model}     & A  \\
  \textbf{Reasoning} & \textless{}think\textgreater{} \newline Okay, let's see. The question is about the final sentence of the passage and which logical error it makes. The options are incomplete comparison, faulty generalization, circular reasoning, fallacy of composition, or moral equivalence. \newline The passage starts by saying that architecture is about building beautifully, and if a structure is just about strength and stability with utilitarian considerations, it's engineering. Then it defines architecture as the art that harmonizes utility and beauty. The last sentence says, "It is the most useful of the fine arts and the noblest of the useful arts." \newline Hmm. So the author is trying to position architecture as both a fine art and a useful art. But wait, the first part says architecture is a fine art (since it's defined as an art that harmonizes utility and beauty), and the second part calls it the noblest of the useful arts. But the problem here might be that they're comparing architecture to two different categories: first, the fine arts, where it's the most useful, and then the useful arts, where it's the noblest. That seems like an incomplete comparison because they're not comparing it to both categories in the same way. The first part is about being the most useful among fine arts, and the second is about being the noblest among useful arts. But the comparison isn't fully explained or justified. Maybe the author is assuming that these two categories are being compared properly, but they aren't. So the error is an incomplete comparison. Let me check the options again. Option a is incomplete comparison. That seems right. The other options don't fit as well. Faulty generalization would be if they made a broad claim without enough evidence. Circular reasoning would be if they defined architecture in terms of itself. Fallacy of composition is assuming what's true for the part is true for the whole. Moral equivalence is comparing two things as morally equal when they aren't. The last sentence isn't doing any of those. So the answer is a. \newline \textless{}/think\textgreater{} \newline a. incomplete comparison  \\
  \midrule
  \multicolumn{2}{@{}l}{\textbf{q030}} \\
  \textbf{Passage}   & Use the following passage to answer question 30. \newline The humorous story is strictly a work of art---high and delicate art---and only an artist can tell it; but no art is necessary in telling the comic and the witty story; anybody can do it. \newline The art of telling a humorous story---understand, I mean by word of mouth, not print---was created in America, and has remained at home. \newline Source: Excerpt from "How to Tell a Story" by Mark Twain    \\
  \textbf{Question}  & 30. Which statement, if it were true, most significantly weakens the argument in the passage? \newline a. Native American storytellers had been telling humorous stories for hundreds of years. \newline b. The plays of Aristophanes are among the earliest examples of humorous stories in writing. \newline c. Mark Twain is widely regarded as one of the wittiest writers in American literature. \newline d. The word art is most commonly applied to visual arts such as painting, sculpture, and drawing. \newline e. It is important to understand certain rules before attempting to tell a comic story.   \\
  \textbf{Gold}      & E       \\
  \textbf{Model}     & A  \\
  \textbf{Reasoning} & \textless{}think\textgreater{} \newline Okay, let's see. The question is asking which statement would weaken the argument in the passage. The passage says that the humorous story is a high art that only an artist can tell, but comic and witty stories don't require art and anyone can do it. Then it mentions that the art of telling humorous stories (by word of mouth, not print) was created in America and has stayed there. \newline So the argument here is that the skill of oral humorous storytelling is uniquely American and requires artistry, while comic and witty stories don't. The question is looking for a statement that would weaken this argument. \newline Looking at the options: \newline a. Native American storytellers had been telling humorous stories for hundreds of years. If this is true, then the idea that the art was created in America might be weakened because Native Americans were already doing it before the concept of "America" as a nation existed. So maybe the art wasn't actually created in America, but existed there earlier. \newline b. The plays of Aristophanes are among the earliest examples of humorous stories in writing. The passage is about oral storytelling, not written. So this might not directly affect the argument since the claim is about the art being created in America through oral means, not written. So maybe not the best choice. \newline c. Mark Twain is considered one of the wittiest writers. The passage is by Twain, and he's saying that witty stories don't need art. If he's a witty writer, that might support his point, not weaken it. So probably not c. \newline d. The word art is commonly applied to visual arts. The passage uses "art" in the context of storytelling. If the term is usually for visual arts, maybe that's a stretch, but the argument is about the existence of the art form in oral storytelling. This might not directly weaken the claim about where it was created. \newline e. It's important to understand rules before telling a comic story. The passage says that comic stories don't require art, but if there are rules to follow, maybe that implies some skill or art? But the original argument says "no art is necessary," so if there are rules, maybe that's a different aspect. Not sure if this weakens the argument about the origin in America. \newline So the main point to weaken is that the art of oral humorous storytelling was created in America. If a. says Native Americans were doing it for centuries, then the art wasn't created in America (as a country), but existed there before. That would weaken the argument. The other options don't directly address the origin or the necessity of art in comic vs humorous stories. So the answer is a. \newline \textless{}/think\textgreater{} \newline a.  \\
  \midrule
  \multicolumn{2}{@{}l}{\textbf{q031}} \\
  \textbf{Passage}   & Use the following passage to answer questions 31 through 33. \newline Whenever we hear the name of Napoleon mentioned, or see it printed in a book, it is usually in connection with a hard-fought victory on the battlefield. He certainly spent most of his life in the camp, and enjoyed the society of soldiers more than that of courtiers. The thunder of guns, the charge of cavalry, and the flash of bayonets as they glittered in the sun, appealed to him with much the same force as music to more ordinary folk... \newline We are apt to forget that this mighty conqueror, whom Carlyle calls "our last great man," had a childhood at all. He was born nearly a century and a half ago, on the 15th August 1769 to be exact, in the little town of Ajaccio, the capital of picturesque Corsica. \newline Source: Excerpt from The Story of Napoleon by Harold F.B. Wheeler    \\
  \textbf{Question}  & 31. As it is used in the context of the sentence, which word best describes the meaning of force (line 9)? \newline a. compel \newline b. intensity \newline c. shove \newline d. pry \newline e. weight   \\
  \textbf{Gold}      & B       \\
  \textbf{Model}     & A  \\
  \textbf{Reasoning} & \textless{}think\textgreater{} \newline Okay, let's see. The question is about the word "force" in the sentence: "the thunder of guns, the charge of cavalry, and the flash of bayonets as they glittered in the sun, appealed to him with much the same force as music to more ordinary folk..." So, I need to figure out what "force" means here. \newline First, the context is that Napoleon was drawn to the sounds and sights of battle, like thunder and cavalry charges, in the same way that others are drawn to music. The word "force" here is being compared to how music affects people. So, "force" probably means a strong attraction or influence. \newline Looking at the options: a. compel, b. intensity, c. shove, d. pry, e. weight. \newline "Compel" means to force someone to do something, which fits if "force" is about making him feel attracted. "Intensity" refers to strength or force in terms of how strong something is. "Shove" is a physical push, which doesn't fit. "Pry" is to look into something, not relevant here. "Weight" is physical mass, also not fitting. \newline The sentence is saying that the battle elements appealed to him with the same force as music does. So "force" here is about the strength or intensity of the appeal. Between a and b. "Compel" is more about causing someone to act, while "intensity" is about the strength of the effect. The original sentence uses "appealed to him with much the same force," so "force" is the strength of the appeal. So "intensity" (b) might be better. But wait, "compel" is also a synonym for force. Hmm. \newline Wait, the comparison is between the force of battle elements and the force of music. So if music compels people, then battle elements compelled Napoleon. But "force" in the sentence is used as a noun. The options are a. compel (verb), but the question is about the noun. Wait, the answer choices are all verbs except for maybe the question. Wait, the question is asking for the meaning of the word "force" as a noun. The options are all verbs. Wait, no, the options are words that could replace "force" in the sentence. The original sentence uses "force" as a noun. So the answer should be a word that can function as a noun and means a strong influence or attraction. \newline Looking again: a. compel (verb), but maybe the noun form is "compulsion"? Not listed. The options are all verbs. Wait, the question is asking for the best word that describes the meaning of "force" in the context. So the answer is the word that, when used in the sentence, would mean the same as "force." The sentence is using "force" to mean a strong influence or attraction. So "compel" is a verb, but if we use it as a noun, maybe "compulsion" is the idea. However, the options are all verbs. Wait, maybe the question is phrased as "which word best describes the meaning of force," so the answer is the verb that would fit in the sentence if "force" were replaced by it. \newline Wait, the original sentence uses "force" as a noun. The options are verbs. So the question is a bit confusing. But maybe the answer is a. compel, because "appealed to him with much the same force as music to more ordinary folk" could be rephrased as "compelled him as music compels others." So "force" here is like the act of compelling. So the answer is a. compel. \newline \textless{}/think\textgreater{} \newline a. compel  \\
  \midrule
  \multicolumn{2}{@{}l}{\textbf{q042}} \\
  \textbf{Passage}   & Use the following passage to answer questions 41 and 42. \newline Science is knowledge; it is what we know. But mere knowledge is not science. For a bit of knowledge to become a part of science, its relation to other bits of knowledge must be found. \newline In botany, for example, bits of knowledge about plants do not make a science of botany. To have a science of botany, we must not only know about leaves, roots, flowers, seeds, etc., but we must know the relations of these parts and of all the parts of a plant to one another. In other words, in science, we must not only know, we must not only have knowledge, but we must \newline know the significance of the knowledge, must know its meaning. This is only another way of saying that we must have knowledge and know its relation to other knowledge. \newline Source: Excerpt from The Science of Human Nature by William Henry Pyle    \\
  \textbf{Question}  & 42. The passage is primarily concerned with \newline a. defining the term knowledge. \newline b. explaining how scientists must use knowledge. \newline c. emphasizing the importance of the whole over the part. \newline d. proving that most people who call themselves scientists are frauds. \newline e. showing how one becomes a scientist.   \\
  \textbf{Gold}      & B       \\
  \textbf{Model}     & C  \\
  \textbf{Reasoning} & \textless{}think\textgreater{} \newline Okay, let's see. The question is asking what the passage is primarily concerned with. The options are a through e. \newline The passage starts by saying science is knowledge, but then it clarifies that just having knowledge isn't enough. It uses botany as an example: knowing about leaves, roots, etc., isn't enough; you need to know how they relate. So the main point seems to be that science isn't just isolated facts but understanding their relationships. \newline Looking at the options: a is about defining knowledge, but the passage is more about the relations between knowledge. b says explaining how scientists must use knowledge. The passage does talk about needing to know relations, which is part of using knowledge in science. c is emphasizing the whole over the part. The example does mention parts of a plant and their relations, so maybe that's it. d is about proving most scientists are frauds, which doesn't fit. e is how to become a scientist, which also doesn't seem right. \newline The key is the emphasis on relations between parts, making the whole. So c might be correct. But the passage is about science needing relations between knowledge bits, not just the whole vs. part. Wait, the question is about the primary concern. The passage is explaining that science requires understanding the relationships between different pieces of knowledge, which is similar to knowing the whole structure rather than just individual parts. So option c, emphasizing the importance of the whole over the part, could be the answer. Alternatively, option b is about how scientists use knowledge, which includes understanding relations. But the passage is more about the necessity of relations for knowledge to be scientific, not just how they use it. Hmm. The example in botany shows that knowing parts isn't enough; you need to know their relations. So the main idea is that science is about the relationships (the whole) rather than just individual parts. So the answer is c. \newline \textless{}/think\textgreater{} \newline C  \\
  \midrule
  \multicolumn{2}{@{}l}{\textbf{q044}} \\
  \textbf{Passage}   & Use the following passage to answer questions 43 and 44. \newline At the end of the nineteenth century the two masters of the stage were Ibsen and Wagner, and both of them were in the show business---Wagner more openly and more frequently than \newline Ibsen. Yet the stern Scandinavian did not disdain to employ an avalanche in When We Dead Awaken, and to introduce a highly pictorial \newline shawl dance for the heroine of his Doll's House. As for Wagner, he was incessant in his search for the spectacular, insisting that the music-drama was the "art-work of the future," since the librettist-composer could call to his aid all the other arts, and could make these arts contribute to the total effect of the opera. He conformed his practise to his principles, and as a result there is scarcely any one of his music-dramas which is not enriched by a most elaborate scenic accompaniment. The forging of the sword, the ride of the Valkyries, the swimming of the singing Rhinemaidens, are only a few of the novel and startling effects which he introduced into his operas; and in his last work, Parsival, the purely spectacular element is at least as ample and as varied as any that can be found in a Parisian fairy-play or in a London Christmas pantomime. \newline Source: Excerpt from A Book about the Theater by Brander Matthews    \\
  \textbf{Question}  & 44. Which sentence from the passage represents an opinion rather than a fact? \newline a. "At the end . . . Ibsen." \newline b. "Yet the . . . House." \newline c. "As for Wagner . . . opera." \newline d. "He conformed . . . accompaniment." \newline e. "The forging . . . pantomime."   \\
  \textbf{Gold}      & E       \\
  \textbf{Model}     & C  \\
  \textbf{Reasoning} & \textless{}think\textgreater{} \newline Okay, let's see. The question is asking which sentence represents an opinion rather than a fact. The options are a through e. \newline First, I need to recall the difference between fact and opinion. A fact is something that can be proven true or false, while an opinion is a personal belief or judgment that may not be verifiable. \newline Looking at the passage, the first sentence (a) states that Ibsen and Wagner were the two masters of the stage at the end of the 19th century. That's a factual claim, though it might be debatable, but it's presented as a statement of historical assessment. \newline Sentence (b) mentions Ibsen using an avalanche and a shawl dance. These are specific elements in his plays, so they are facts. \newline Sentence (c) says Wagner was "incessant in his search for the spectacular" and that he considered music-drama the "art-work of the future." The term "incessant" is subjective, indicating a strong, continuous effort, which is an opinion. Also, his belief about the art-work of the future is his own opinion, not a fact. \newline Sentence (d) talks about Wagner conforming his practice to his principles, leading to elaborate scenic accompaniment. This is a factual statement about his work practices. \newline Sentence (e) lists specific effects like the forging of the sword, Valkyries' ride, etc., which are elements from his operas, so facts. \newline So the opinion is in sentence (c). \newline \textless{}/think\textgreater{} \newline c. "As for Wagner . . . opera."  \\
  \midrule
  \bottomrule
\end{longtable}